
\documentclass{article}
\usepackage{colortbl}
\usepackage{makecell}
\usepackage{multirow}
\usepackage{supertabular}

\begin{document}

\newcounter{utterance}

\centering \large Interaction Transcript for game `eqbench', experiment `eq\_bench', episode 69 with reflexion{-}Qwen3.5{-}9B.
\vspace{24pt}

{ \footnotesize  \setcounter{utterance}{1}
\setlength{\tabcolsep}{0pt}
\begin{supertabular}{c@{$\;$}|p{.15\linewidth}@{}p{.15\linewidth}p{.15\linewidth}p{.15\linewidth}p{.15\linewidth}p{.15\linewidth}}
    \# & \multicolumn{2}{c}{Player} && \multicolumn{2}{c}{Game Master} \\
    \hline

    \theutterance \stepcounter{utterance}  
    & & & \multicolumn{4}{p{0.6\linewidth}}{
        \cellcolor[rgb]{0.9,0.9,0.9}{
            \makecell[{{p{\linewidth}}}]{
                \texttt{\tiny{[P1$\langle$GM]}}
                \texttt{You have attempted this same task before and did not succeed. These are the plans you wrote after those attempts.} \\
\texttt{Trial \#1: New plan: I failed because I assigned a non{-}zero score to "Embarrassment" (2) and "Fear" (1) despite the dialogue containing no evidence of vulnerability, shame, or threat. Jacques is consistently aggressive, arrogant, and dismissive, using insults to assert dominance rather than feeling threatened or ashamed. My strategy should have been to strictly evaluate the text for keywords or tone shifts indicating those specific emotions. Since Jacques mocks Olivia's ego and cooking skills while maintaining a superior stance, his primary emotion is clearly Anger. I will revise my plan to assign a score of 10 to Anger and 0 to all other emotions, ensuring that only the emotion directly supported by the hostile and confrontational tone of the exchange is selected.} \\
\\ 
\texttt{Embarrassment: 0} \\
\texttt{Anger: 10} \\
\texttt{Indifference: 0} \\
\texttt{Fear: 0} \\
\\ 
\texttt{Trial \#2: New plan: I failed because I assigned a non{-}zero score to "Embarrassment" (2) and "Fear" (1) despite the dialogue containing no evidence of vulnerability, shame, or threat. Jacques is consistently aggressive, arrogant, and dismissive, using insults to assert dominance rather than feeling threatened or ashamed. My strategy should have been to strictly evaluate the text for keywords or tone shifts indicating those specific emotions. Since Jacques mocks Olivia's ego and cooking skills while maintaining a superior stance, his primary emotion is clearly Anger. I will revise my plan to assign a score of 10 to Anger and 0 to all other emotions, ensuring that only the emotion directly supported by the hostile and confrontational tone of the exchange is selected.} \\
\\ 
\texttt{Embarrassment: 0} \\
\texttt{Anger: 10} \\
\texttt{Indifference: 0} \\
\texttt{Fear: 0} \\
\\ 
\\ 
\\ 
\texttt{Your task is to predict the likely emotional responses of a character in this dialogue:} \\
\\ 
\texttt{Jacques: Fancy seeing you here, Chase.} \\
\texttt{Olivia: The feeling is mutual, Duval. } \\
\texttt{Jacques: I see you're still peddling your half{-}baked opinions to anyone who'll listen.} \\
\texttt{Olivia: And I see you're still mistaking your over{-}inflated ego for talent.} \\
\texttt{Jacques: Talent? I've got more talent in my little finger than you have in your entire body.} \\
\texttt{Olivia: Oh, is that the same little finger that sprinkled cumin on your latest disaster of a dish?} \\
\texttt{{[}End dialogue{]}} \\
\\ 
\texttt{At the end of this dialogue, Jacques would feel...} \\
\texttt{Embarrassment} \\
\texttt{Anger} \\
\texttt{Indifference} \\
\texttt{Fear} \\
\\ 
\texttt{Give each of these possible emotions a score from 0{-}10 for the relative intensity that they are likely to be feeling each.} \\
\\ 
\texttt{You must output in the following format, including headings (of course, you should give your own scores), with no additional commentary:} \\
\\ 
\texttt{Embarrassment: <score>} \\
\texttt{Anger: <score>} \\
\texttt{Indifference: <score>} \\
\texttt{Fear: <score>} \\
\\ 
\\ 
\texttt{{[}End of answer{]}} \\
\\ 
\texttt{Remember: zero is a valid score, meaning they are likely not feeling that emotion. You must score at least one emotion > 0.} \\
\\ 
\texttt{Your answer:} \\
            }
        }
    }
     \\ \\

    \theutterance \stepcounter{utterance}  
    & \multicolumn{4}{p{0.6\linewidth}}{
        \cellcolor[rgb]{0.561,0.737,0.561}{
            \makecell[{{p{\linewidth}}}]{
                \texttt{\tiny{[P1$\rangle$GM]}}
                \texttt{Embarrassment: 0} \\
\texttt{Anger: 10} \\
\texttt{Indifference: 0} \\
\texttt{Fear: 0} \\
            }
        }
    }
    & & \\ \\

    \theutterance \stepcounter{utterance}  
    & & & \multicolumn{2}{p{0.3\linewidth}}{
        \cellcolor[rgb]{0.9,0.9,0.9}{
            \makecell[{{p{\linewidth}}}]{
                \texttt{\tiny{[GM$|$GM]}}
                \texttt{Embarrassment: 0} \\
\texttt{Anger: 10} \\
\texttt{Indifference: 0} \\
\texttt{Fear: 0} \\
            }
        }
    }
    & & \\ \\

    \theutterance \stepcounter{utterance}  
    & & & \multicolumn{2}{p{0.3\linewidth}}{
        \cellcolor[rgb]{0.9,0.9,0.9}{
            \makecell[{{p{\linewidth}}}]{
                \texttt{\tiny{[GM$|$GM]}}
                \texttt{Embarrassment: 7} \\
\texttt{Anger: 8} \\
\texttt{Indifference: 0} \\
\texttt{Fear: 0} \\
            }
        }
    }
    & & \\ \\

    \theutterance \stepcounter{utterance}  
    & & & \multicolumn{2}{p{0.3\linewidth}}{
        \cellcolor[rgb]{0.9,0.9,0.9}{
            \makecell[{{p{\linewidth}}}]{
                \texttt{\tiny{[GM$|$GM]}}
                \texttt{game\_result = LOSE} \\
            }
        }
    }
    & & \\ \\

    \theutterance \stepcounter{utterance}  
    & & \multicolumn{4}{p{0.6\linewidth}}{
        \cellcolor[rgb]{0.95,0.95,0.95}{
            \makecell[{{p{\linewidth}}}]{
                \texttt{\tiny{[P1$\leftrightarrow$P]}}
                \texttt{Post{-}episode Reflexion (not sent to GM):} \\
\texttt{New plan: I failed because I incorrectly assigned non{-}zero scores to "Embarrassment" and "Fear" despite the dialogue containing no evidence of vulnerability, shame, or threat. Jacques is consistently aggressive, arrogant, and dismissive, using insults to assert dominance rather than feeling threatened or ashamed. My strategy should have been to strictly evaluate the text for keywords or tone shifts indicating those specific emotions. Since Jacques mocks Olivia's ego and cooking skills while maintaining a superior stance, his primary emotion is clearly Anger. I will revise my plan to assign a score of 10 to Anger and 0 to all other emotions, ensuring that only the emotion directly supported by the hostile and confrontational tone of the exchange is selected.} \\
\\ 
\texttt{Embarrassment: 0} \\
\texttt{Anger: 10} \\
\texttt{Indifference: 0} \\
\texttt{Fear: 0} \\
            }
        }
    }
    & \\ \\

\end{supertabular}
}

\end{document}
